\documentclass[11pt]{article}
\input{preamble}
\input{macros}
\addbibresource{refs.bib}
\linespread{1.11}

\title{Categorical Message Passing Language (CaMPL)}
\author{
    Daniel Hashimoto\thanks{Universidade Federal do Rio de Janeiro} \and Alexanna Little\thanks{University of Calgary} \and \and Priyaa Varshinee Srinivasan\thanks{Tallinn University of Technology. The author was funded by...}}
%\def\authorrunning{D.\ I.\ Spivak \& P.\ V.\ Srinivasan}
\date{\today}

\begin{document}

\maketitle

\begin{abstract}
These notes are a compact introduction to  a novel concurrent programming language called the 
categorical message passing language built on the mathematics of linear actegories. 
Categorical Message Passing Language, to the best of our knowledge, 
the only\footnote{Is this true?} concurrent programming language with mathematical underpinnings. 
It is a functional-style programming language with the abstract machine implemented in Haskell. 
The current version of CaMPL compiler is to be considered a proto-alpha, a proof-of-concept 
of a programming language built from the mathematics of linear actegories. \textcolor{blue}{Hashimoto's comment:Does this have to be on the abstract?
(It repeats "linear actagories")
I would rather put something cool like:
no deadlocks,
type system made for designing protocols,
message passing (erlang-like actors),
... .
Here is a good place to put your selling points.}
\end{abstract}

%\addcontentsline{toc}{section}{References}

%\tableofcontents

\section{Introduction}

A paragraph on why CaMPL (Curry-Howard-Lambek-like correspondence). 

In this article, we have adopted a narrative-style of writing instead of a programming manual for CaMPL. To make it easy for the reader, the main CaMPL ideas of each section are summarized as {\bf Keywords:} at the beginning of the section.

\note{Xanna note:
i am assuming introduction will establish justification for why we have submitted this paper to this venue.
and that we will mention things like there is a beautiful type system and that we are doing message passing and that there is a categorical semantics}

\subsection{A brief overview of CaMPL}

\colorbox{green!40}{Keywords:} Processes, protocols and co-protocols, and channels.

\vspace{0.25em}

As is tradition, we will introduce CaMPL with a ``Hello World!'' example program. The ``Hello World!'' example demonstrates communication between a process and the outside world using the {\tt Console} channel type. 

The following code snippet shows a CaMPL process named {\tt helloworld} printing {\tt ``Hello World!''} by sending the string as a message on a channel named {\tt console}. 
The {\tt console} channel has the type {\tt Console} which is a special channel type hard-coded in the compiler to connect the program to the terminal from which it is run.
\begin{lstlisting}
proc helloworld :: | Console => =
	| console => -> do
		hput ConsolePut on console
		put "Hello World!" on console    -- sends message to the console
		hput ConsoleClose on console
		halt console                     -- closes console channel and halts
\end{lstlisting}
The \tc{process} \lstinline$helloworld$ communicates to an built-in printing process via a channel of type \lstinline*Console*. In CaMPL, concurrent types are \tc{(co)protocols}, the instantiations of which are called as \tc{channels}. The \tc{process command} \lstinline$hput$ may seem strange at the first sight --- read as `handle put',  it initiates a channel for a  communication session. \tc{Protocol handles} are analogous to session types. As indicated by their names, the \lstinline$put$ process command sends a message, and the \lstinline$halt$ process command halts the calling process.

The {\tt helloworld} process can be invoked by calling it in the main process, {\tt run}, and giving it the {\tt console} channel as follows:
\begin{lstlisting}
 proc run = 
	| console => -> 
		helloworld( | console => )  -- creates helloworld process
\end{lstlisting}

The channel to the console can only be created by the main process.

\section{Designing a simple Echo server}
% Describe in-built channel types - Put, Get, TopBot
% Describe process command - plug
% Describe channel polarities and their roles
% A drop of mathematics and logic

\knote{Commands \lstinline$plug$, \lstinline$put$, \lstinline$get$. Types \lstinline$Put$, \lstinline$Get$, \lstinline$TopBot$.} 

\concept{Connecting processes via a channel, one-shot send and receive, closing communication.}

As we become more familiar with CaMPL, let us begin by designing a simple echo server in which a client and a server will exchange one message each. 

In order to arrange such a communication to proceed, we need to connect the client and the server, which is facilitated by the \lstinline$plug$ process command. 
\begin{lstlisting}
proc run = 
	| => -> 
		plug 
			client( | => ch )
			server( | ch => )
\end{lstlisting}

The above {\tt run} process, in Line 3, creates a {\tt client} and a {\tt server} process, and plugs them together via a named {\tt ch}. The CaMPL compiler has type-inference ability and will resolve the type of the channel {\tt ch}. 

It may be noticed that, in line 4, that {\tt ch} is placed after "\lstinline$=>$" for the client, and in line 5, that, {\tt ch} is placed before "\lstinline$=>$" for the server. Whether the channel {\tt ch} occurs before or after the arrow determines its \tc{channel polarity} --- the different polarities signifies different ends of the channel. In this case, {\tt ch} is of \tc{output polarity}, also written as \ttc{(+)}, for the {\tt client}. Dually, {\tt ch} is of \tc{input polarity}, also written as \ttc{(-)}, for the {\tt server}. 

Note: The choice of terminology -- input and output -- polarities is best understood by imagining a process as a box, with wires entering it and wires leaving it (with respect to the direction of gravity). We call incoming wires as input polarity, and outgoing as output polarity. 
\comment{from Hashimoto: Figure, please.}

Thus, it must be obvious now that CaMPL allows processes to be plugged together only along channels of same type and opposite polarities. 

Let us now to code the client process:
\begin{figure}[h]
\begin{lstlisting}
proc client :: Int | => Put(Int | Get( [Char] | TopBot)) =
	val | => ch -> do
		put val on ch    -- send a message
		get msg on ch    -- receive a message 
		halt ch          -- close channel and halt process
\end{lstlisting}
\caption{Basic client process}
\label{Fig: Basic client}
\end{figure}

% Explain type versus polarity.
\lstinline$Put$, \lstinline$Get$, and \lstinline$TopBot$ are in-built concurrent types for channels. The \lstinline$Put$ and \lstinline$Get$ types is used to exchange messages, while the \lstinline$TopBot$ type is used to end communication.
\begin{itemize}
\item Line 1 of \lstinline$proc$ {\tt client} clarifies that the {\tt client} process has an output polarity channel of compound type \lstinline$Put(Int | Get( [Char] | TopBot))$. 
\item Line 2 assigns variable name {\tt ch} to the output polarity channel. 
\item Line 4 sends the integer value of {\tt val} over {\tt ch} which is of type \lstinline$Put(Int | Get( [Char] | TopBot))$. On completion, the {\tt ch} proceeds with type \lstinline$Get( [Char] | TopBot)$.
\item Line 5 receives a string over {\tt ch} which is now of type \lstinline$Get( [Char] | TopBot)$, and binds the string value to the variable {\tt msg}. On completion, the {\tt ch} proceeds with type \lstinline$TopBot)$.
\item Line 6 closes {\tt ch} which is now of type {\tt TopBot} (hence, no more communication over {\tt ch}) and halts the {\tt client} process.
\end{itemize}

It will become clear in a moment the role played by polarities in communication, when we code the server process. But, for now, it must be clear that a process may both send and receive over an output polarity channel.

Let us design the server process:
\begin{figure}[h]
\begin{lstlisting}
proc server :: [Char] | Put(Int | Get( [Char] | TopBot)) => =
	msg | ch => -> do
		get inval on ch     -- receive message
		put msg on ch       -- send message 
		halt ch             -- close channel and halt process
\end{lstlisting}
\caption{Basic server process}
\label{Fig: Basic server}
\end{figure}

\begin{itemize}
\item Line 1 gives the type signature of \lstinline$proc$ {\tt server}. The {\tt server} process has an input polarity channel of compound type \lstinline$Put(Int | Get( [Char] | TopBot))$. 
\item Line 2 assigns variable name {\tt ch} to the input polarity channel. 
\item Line 3 receives an integer value over channel {\tt ch} which is of type \lstinline$Put(Int | Get( [Char] | TopBot))$, and binds it to the variable {\tt inval}. On completion, the {\tt ch} proceeds with type \lstinline$Get( [Char] | TopBot)$.
\item Line 4 sends a string over channel {\tt ch} which is now of type \lstinline$Get( [Char] | TopBot)$. On completion, the {\tt ch} proceeds with type \lstinline$TopBot)$. 
\end{itemize}

Note that, when the process command `\lstinline$put val on ch$' of the {\tt client} process (Line 4) is matched by the process command `\lstinline$get inval on ch$' in the {\tt server} process. This command pair allows communication of a single message. We call the pair {\tt (put, get)} as a \tc{complementary pair} -- meaning the commands are used in pairs one at each end of the same channel. This arises a question, given a complementary pair, which command should be used at which end of the channel? This is where polarities come into play.

Given a channel of type \lstinline$Put$, the process which has the {\bf output polarity} end may \lstinline$put$ a message on the channel, and correspondingly, the process which has the {\bf input polarity} end may \lstinline$get$ a message over the channel. The roles cannot be switched. Given a channel of type \lstinline$Get$, the process which has the {\bf input polarity} end may \lstinline$put$ a message on the channel, and correspondingly, the process which has the {\bf output polarity} end may \lstinline$get$ a message over the channel, and not vice versa. Thus, given either a \lstinline$Put$ or a \lstinline$Get$ type, the polarities allow unambiguous distinction of the role of a process (if it is a sender or a receiver). 

\comment{Hashimoto's suggestion: Lead into the next section with type-inferencing.}

\note{@Xanna, I think you are best suited to write the following paragraph :)}
The notion polarities arise from the underlying mathematics of linear actegories. The linear actegories are equipped with an adjunction parametrized by the sequential datatypes. The following four sequent rules encode the usage of {\tt(put,get)} process commands over \lstinline$Put$ and \lstinline$Get$ types.

\section{Concurrent type system: Handling multiple sequential messages}

\colorbox{green!40}{Keywords:} {\tt protocol} keyword 

\noindent \colorbox{green!40}{Concepts:} Coding custom channel types using protocols.

The server and client processes discussed in the last section are capable to sending one message and receiving one message. For example, the client process begins with an output polarity channel of type \lstinline$Put(Int | Get( [Char] | TopBot))$ which allows a message to be sent over the channel. Once the message is sent, the channel proceeds with type \lstinline$Get( [Char] | TopBot)$ allowing the process to receive a message. Once a message is received, the channel proceeds with type \lstinline$TopBot$ allowing it to be closed. Let us equip the client process to perform arbitrary number of send-receive pairs until it {\em wants to} close the communication channel.

In CaMPL this is achieved using (co)protocols. For simplicity purposes we will focus on protocols. 

The following code snippet shows a protocol which the client and the server processes may use for communicating arbitrary number of messages with one another.

\begin{figure}[h]
\begin{lstlisting}
protocol Echo => S =
    EchoSend :: Put(Int | Get( [Char] | S)) => S
    EchoClose :: TopBot => S
\end{lstlisting}
\caption{Protocol Echo}
\label{Fig: Protocol Echo}
\end{figure}

\begin{itemize}
	\item Line 1 declares a \lstinline$protocol$ named \lstinline$Echo$ that has an output concurrent type \lstinline$S$, and no input sequential and concurrent types.
	\item Line 2 declares a protocol handle named \lstinline$EchoSend$ which when {\em placed} on a channel of type \lstinline$Echo$, allows a process to do  \lstinline$put$-\lstinline$get$, or  \lstinline$get$-\lstinline$put$ over it depending on the polarity. 
	
	Compare the types in the Line 2 of the above code  and line 1 in Figure \ref{Fig: Basic client}. With the \lstinline$EchoSend$ handle, the channel will proceed with type \lstinline$S$ once a send-receive operation is completed. Thereby, the same handle can be placed once more on the channel to perform another round of send-receive operation.
	
	\item Line 3 declares a protocol handle named \lstinline$EchoClose$ which when {\em placed} on a channel of type \lstinline$Echo$, allows a process to end communication using \lstinline$close$ or \lstinline$halt$ commands. 
\end{itemize}

Protocols handles are analogous to session types. They represent a valid sequence of process commands that may be performed on a channel.

More to come.. How to use this protocol with client, server and run process.

\section{Split and Fork: Asynchronous querying and response}

% This is a good place to talk about deadlocks.

\section{Races: Handling multiple concurrent queries}

\section{Discussion and future work}

\end{document}
